\documentclass[11pt]{article}
\usepackage[utf8]{inputenc}
\usepackage[T1]{fontenc}
\usepackage[margin=1in]{geometry}
\usepackage{booktabs}
\usepackage{array}
\usepackage{amsmath}
\usepackage{amssymb}
\usepackage{hyperref}
\usepackage{enumitem}
\usepackage{xcolor}
\usepackage{parskip}

\title{Discussion Log \\[0.3em] \large StockVolatilitySight --- CPSC 440}
\author{}
\date{}

\begin{document}
\maketitle

Running record of important decisions, divergences from the proposal, and rationale. Keep entries concise and focused --- this is a working document for the team, not a plan artifact.

\paragraph{Conventions.}
\begin{itemize}[leftmargin=*]
  \item The proposal (\texttt{plan/Research 440 Proposal.md}) frames the research question and is the source of truth for \emph{intent}.
  \item The plan doc (\texttt{plan/hmm-lstm\_volatility\_forecast\_11d292d8.plan.md}) is a hypothesis / implementation log. Empirical results can and do diverge from what the plan predicted --- that's the research finding, not a bug. Log divergences here with neutral framing.
  \item When implementation diverges from the \emph{proposal}, log it here with a rationale.
\end{itemize}

\hrulefill

\section*{Current State (updated 2026-04-23, post-rerun)}

\paragraph{Pipeline status.} Phases 1--6 of the plan are complete. Full pipeline re-run from a fresh clone today (nb 01 $\to$ 04 $\to$ 05 $\to$ \texttt{src/ensemble.py} $\to$ nb 06). HMM selected (GMMHMM, 2 states, 2 mixture components, full covariance). All three LSTMs trained. Ensemble evaluated on 2020-06-30 $\to$ 2026-03-16 ($n=1{,}434$).

\paragraph{Headline empirical result (fresh run, \texttt{notebooks/06\_ensemble\_eval.ipynb}):}

\begin{center}
\begin{tabular}{lrrr}
\toprule
Model & MSE & MAPE & Prior-run MSE \\
\midrule
Naive (\texttt{rolling\_std\_21}) & 2.2e-5 & 34.6\% & 2.2e-5 \\
\textbf{Baseline LSTM}            & \textbf{1.6e-5} & \textbf{25.1\%} & 1.5e-5 \\
Calm LSTM                         & 4.1e-5 & 34.9\% & 2.0e-5 \\
Volatile LSTM                     & 8.5e-5 & 119.4\% & 8.5e-5 \\
Ensemble                          & 2.6e-5 & 31.9\% & 1.7e-5 \\
\bottomrule
\end{tabular}
\end{center}

\paragraph{Diebold--Mariano significance tests} (nb 06 \S6, first time actually executed):

\begin{center}
\renewcommand{\arraystretch}{1.1}
\begin{tabular}{@{}lllrrl@{}}
\toprule
Comparison & Loss & DM & $p$ & Verdict ($\alpha=.05$) \\
\midrule
Naive vs Baseline (full)               & MSE & $+1.36$ & 0.173 & tie \\
Naive vs Baseline (full)               & MAE & $+2.46$ & 0.014 & \textbf{Baseline wins} \\
Baseline vs Ensemble (full)            & MSE & $-1.75$ & 0.080 & tie (marginal, $\alpha=.10$) \\
Baseline vs Ensemble (full)            & MAE & $-2.00$ & 0.046 & \textbf{Baseline wins} \\
Baseline vs Ensemble (volatile, $n=27$)& MSE & $+0.62$ & 0.542 & tie \\
Baseline vs Ensemble (volatile, $n=27$)& MAE & $+0.22$ & 0.831 & tie \\
\bottomrule
\end{tabular}
\end{center}

On aggregate the baseline LSTM beats the HMM-LSTM ensemble --- significantly on MAE ($p=0.046$), marginally on MSE ($p=0.080$). On the 27 strictly-volatile test days the ensemble is \emph{directionally} better (MSE 5e-5 vs 8e-5) but the sample is too small to reject ($p > 0.5$). Regime separation helps where it should; the calm regime dominates aggregate metrics.

\paragraph{Training stochasticity concern.} This run's ensemble MSE is 2.6e-5 vs the 1.7e-5 in the previously-saved notebook output --- a $\sim$50\% difference from the same config, with different Optuna/PyTorch numerical noise. The calm LSTM doubled its MSE (4.1e-5 vs 2.0e-5). \textbf{Single-seed results are not reliable for ranking claims} --- before locking in any narrative we should run $k$ seeds and report mean$\pm$std.

\paragraph{Known reproducibility gaps (four now).}
\begin{itemize}[leftmargin=*]
  \item \texttt{data/processed/train|val|test.parquet} predate the Phase-6 addition of \texttt{relative\_volume\_21d}. Cloning and running nb 04 directly $\to$ \texttt{KeyError}. Fix: re-run nb 01 (which calls \texttt{features.build\_features} and saves with the column).
  \item \textbf{\texttt{src/ensemble.py} is an undocumented step between nb 05 and nb 06.} Nb 06 cell 4 reads \texttt{data/processed/test\_predictions.parquet}, which \texttt{ensemble.py} produces. Neither the README nor the notebooks call this out. The nb 05 $\to$ nb 06 chain fails with \texttt{FileNotFoundError} without this step. Discovered today when running end-to-end.
  \item \texttt{data/processed/test\_predictions.parquet} is not tracked (produced by the step above).
  \item LSTM checkpoints (\texttt{models/*.pt}) are gitignored.
\end{itemize}

\paragraph{Open items to decide.}
\begin{enumerate}[leftmargin=*]
  \item Is the aggregate-loss / volatile-day-win tradeoff acceptable for the writeup, or do we want to revisit gating? (Proposal explicitly listed a learned gating network as out of scope.)
  \item Volatile LSTM outputs ${\sim}0.0175$ nearly constantly and pulls the ensemble up on calm days --- does this need a fix (clipping, higher gating threshold) or is it part of the finding?
  \item \texttt{config.VAL\_END = "2020-04-30"} means val now includes COVID. The plan text still says ``2016--2019.'' Decide whether to update the plan text or realign the config.
  \item \textbf{Multi-seed robustness study.} Re-run LSTM training under $k \geq 5$ random seeds to get mean$\pm$std of the metrics. The current single-run swing from ensemble-MSE 1.7e-5 to 2.6e-5 means any ``baseline beats ensemble'' claim is fragile.
  \item Document the nb 05 $\to$ \texttt{ensemble.py} $\to$ nb 06 step in \texttt{README.md} (or bundle \texttt{ensemble.py} invocation into nb 06's setup cell).
\end{enumerate}

\hrulefill

\section*{2026-04-23 --- Plan summary and proposal/plan discrepancies}

\subsection*{Plan summary}

Regime-aware hybrid model for 21-day forward realized volatility on SPY:

\begin{enumerate}[leftmargin=*]
  \item \textbf{Data + features:} OHLCV + put/call + AAII sentiment $\to$ log returns, abs returns, ranges, rolling stats; chronological split; normality diagnostics to justify HMM choice.
  \item \textbf{HMM regime detection:} Gaussian vs GMM-HMM, 1st- vs 2nd-order (via augmented state space), BIC/LL comparison, Viterbi sanity check against 2008/2020/2022 crises.
  \item \textbf{LSTMs:} baseline (full data) + calm + volatile, PyTorch CPU, log-target with MSE, Optuna TPE tuning.
  \item \textbf{Ensemble:} $y = P(\text{calm} \mid O) \cdot y_{\text{calm}} + P(\text{volatile} \mid O) \cdot y_{\text{volatile}}$.
  \item \textbf{Evaluation:} MSE/RMSE/MAE vs baseline LSTM on 2020-present.
\end{enumerate}

\textbf{Plan-era result (pre-Phase-6):} ensemble MSE 0.00167 vs baseline 0.00246 on the 2020-present holdout --- the regime-split hypothesis was initially validated. See ``Current State'' above for the post-Phase-6 empirical update.

\subsection*{Discrepancies between proposal and plan}

\begin{center}
\renewcommand{\arraystretch}{1.15}
\begin{tabular}{@{}c p{3.4cm} p{4.4cm} p{5.8cm}@{}}
\toprule
\# & Proposal says & Plan / code does & Rationale \\
\midrule
1 & Targets SPY, XEQT, and MAG7 individual stocks & SPY only & XEQT has insufficient history; MAG7 dropped from scope. Documented in plan. \\
2 & Train 2005--2015, Val 2015--2020, Test remaining & Train 2004--2015, Val 2016--2019, Test 2020-present & Not explicitly rationalized in plan (git log shows a revert on the split). Non-overlapping boundaries are cleaner but the drift from the proposal is undocumented. \textbf{Open item.} \\
3 & Inputs include raw OHLCV & Phase 6 removed raw OHLCV; kept \texttt{log\_return}, \texttt{abs\_return}, \texttt{oc\_return}, \texttt{intraday\_range}, \texttt{relative\_volume\_21d} & Raw prices caused non-stationary drift in predictions; stationarity fix dropped MAPE from ${\sim}600$\% to ${\sim}25$\%. Tried, didn't work --- legitimate divergence. \\
4 & Compare Gaussian vs GMM-HMM; explore higher-order Markov ``if time permits'' & Both done (GMM-HMM and 2nd-order via augmented state space) & No discrepancy. \\
5 & Reference baselines: baseline LSTM + ``considering'' a simple tree/NN & Baseline LSTM + na\"ive \texttt{rolling\_std\_21} (no tree model) & Plan picked a different simple baseline. Proposal's tree regressor was never implemented. \\
6 & Regime-specific LSTMs train on regime-partitioned data & Volatile-LSTM val fallback: 0 volatile-dominant windows in 2016--2019, so early-stopping uses the unfiltered val set & Practical necessity --- the 2016--2019 val window is almost entirely calm. Methodological caveat: volatile LSTM's model selection isn't purely regime-matched. \\
7 & Learned gating network --- considered, out of scope & Out of scope & No discrepancy. \\
8 & Non-price inputs: put/call volume (Ho et al., 2012) and AAII weekly sentiment (Gupta et al., 2023) as key features & \textbf{Scaffolded but inert.} \texttt{SENTIMENT\_FEATURES} + \texttt{PUTCALL\_FEATURES} defined in \texttt{config.py} since commit 888a861 (2026-04-16), loaders exist in \texttt{src/data\_loader.py}, but the CSVs were never downloaded to \texttt{data/raw/} and neither group is referenced by \texttt{HMM\_FEATURES} or \texttt{LSTM\_BASELINE\_FEATURES}. Verified via git log: never added to the active feature lists, never removed --- just never activated. & HMM + LSTMs train on price/volume only. Significant divergence from the proposal's core ``non-price inputs matter for vol'' thesis. Options for the paper: (a) download the CSVs and retrain to restore proposal alignment; (b) explicitly note the scope-down in methods + limitations. \\
\bottomrule
\end{tabular}
\end{center}

\textbf{Most consequential to flag in the writeup:} \#1 (scope narrowing to SPY), \#3 (stationarity-driven feature change), \#6 (volatile-LSTM val fallback), \textbf{\#8 (sentiment + put/call never wired in --- materially weakens the ``non-price signals matter'' thesis)}. \#2 is the open item --- worth deciding whether to retro-justify or realign.

\hrulefill

\section*{2026-04-23 --- Repo survey: code, notebooks, data sanity}

\subsection*{Data we have}

Checked-in processed data (\texttt{data/processed/}):

\begin{itemize}[leftmargin=*]
  \item \texttt{features\_all.parquet} --- full feature table, 5{,}584 rows, 2004-01-05 $\to$ 2026-03-16.
  \item \texttt{train.parquet} (3{,}020 rows, 2004-01-05 $\to$ 2015-12-31), \texttt{val.parquet} (1{,}089 rows, 2016-01-04 $\to$ 2020-04-30), \texttt{test.parquet} (1{,}475 rows, 2020-05-01 $\to$ 2026-03-16).
  \item \texttt{regime\_probabilities.parquet} --- HMM soft probs + Viterbi labels, 5{,}564 rows. Overall regime mix: 96.3\% calm, 3.7\% volatile (matches plan).
\end{itemize}

Columns present in train/val/test (20): \texttt{Open}, \texttt{High}, \texttt{Low}, \texttt{Close}, \texttt{Volume}, \texttt{log\_return}, \texttt{abs\_return}, \texttt{oc\_return}, \texttt{intraday\_range}, \texttt{log\_volume}, \texttt{rolling\_mean\_\{5,10,21\}}, \texttt{rolling\_std\_\{5,10,21\}}, \texttt{rolling\_abs\_mean\_\{5,10,21\}}, \texttt{realized\_vol\_21d}.

Sanity spot-checks on train pass: \texttt{log\_return} mean ${\approx}\,0.0004$, std ${\approx}\,0.012$ (\checkmark\ daily SPY scale); \texttt{log\_volume} $\in [16.5, 20.6]$ (\checkmark\ log of millions); target \texttt{realized\_vol\_21d} positive, same scale as \texttt{rolling\_std\_21}; no NaNs after rolling warmup. Target formula in \texttt{features.py:154} correctly implements the proposal: forward-looking, demeaned, window $=21$ ---
\[
\sigma_t \;=\; \sqrt{\frac{1}{m}\sum_{j=1}^{m} (r_{t+j} - \bar{r})^2}, \qquad m=21.
\]

\subsection*{Code layout}

\texttt{src/} holds all modules (not \texttt{models/}, which only holds serialized \texttt{.joblib} artifacts). Key files:

\begin{itemize}[leftmargin=*]
  \item \texttt{features.py} --- engineers all features and target; matches proposal math.
  \item \texttt{hmm\_model.py} --- Gaussian vs GMM-HMM comparison, penalised EM with restarts, 1st- vs 2nd-order via state-space augmentation.
  \item \texttt{lstm\_model.py} --- stacked LSTM + linear head, sliding-window \texttt{Dataset}, log-space prediction.
  \item \texttt{train\_LSTM\_baseline.py}, \texttt{train\_LSTM\_regime.py} --- Optuna TPE tuning $\to$ final retrain.
  \item \texttt{ensemble.py} --- weighted-sum ensemble keyed on HMM soft probs.
  \item \texttt{config.py} --- single source of truth for features, splits, HMM/LSTM search spaces.
\end{itemize}

HMM artifacts in \texttt{models/}: \texttt{hmm\_winner.joblib} $=$ GMMHMM(states=2, n\_mix=2, cov=full), val LL $= 9057.28$, BIC $= -15925.74$. Volatile state is state 1. \texttt{HMM\_DURATION\_PENALTY=3000} (approximately the number of training observations) prevents degenerate rapid-switching during EM restarts.

\subsection*{Current test-set results}

From \texttt{notebooks/06\_ensemble\_eval.ipynb}, 2020-06-30 $\to$ 2026-03-16, $n=1{,}434$.

\begin{center}
\begin{tabular}{lrrrr}
\toprule
Model & MSE & RMSE & MAE & MAPE \\
\midrule
Naive (\texttt{rolling\_std\_21}) & 0.000022 & 0.00470 & 0.00313 & 34.6\% \\
\textbf{Baseline LSTM}            & \textbf{0.000015} & \textbf{0.00390} & \textbf{0.00246} & \textbf{25.3\%} \\
Calm LSTM                         & 0.000020 & 0.00442 & 0.00279 & 29.5\% \\
Volatile LSTM                     & 0.000085 & 0.00922 & 0.00857 & 119.4\% \\
Ensemble                          & 0.000017 & 0.00415 & 0.00271 & 28.8\% \\
\bottomrule
\end{tabular}
\end{center}

Regime-stratified (HMM confidence $> 0.8$): on strictly-calm days ($n=1{,}407$) baseline wins (MSE 1e-5 vs 2e-5); on strictly-volatile days ($n=27$) \textbf{ensemble wins} (MSE 5e-5 vs 1e-4).

\subsection*{Findings worth flagging}

\begin{enumerate}[leftmargin=*]
  \item \textbf{Empirical ordering differs from the plan's hypothesis.} The plan doc anticipated ensemble $>$ baseline (MSE 0.00167 vs 0.00246, pre-Phase-6). Current post-Phase-6 notebook-06 output shows baseline (1.5e-5) slightly ahead of ensemble (1.7e-5) on aggregate --- but ensemble beats baseline on strictly-volatile days ($n=27$, MSE 5e-5 vs 1e-4). Not a failure: the plan is a hypothesis doc, and this \emph{is} the research result. Worth framing clearly in the writeup as ``regime separation helps on volatile regimes specifically; the calm regime dominates 2020--2026 so the baseline is hard to beat on aggregate.''

  \item \textbf{Volatile LSTM is overpredicting on calm days.} Its test MAPE is 119\% --- when the HMM is wrong about the regime (or uncertain), the volatile model's near-constant high output (${\sim}0.0175$) pulls the ensemble up. Options: (a) clip volatile predictions when \texttt{p\_volatile} is small, (b) use a higher gating threshold, (c) revisit the learned-gating-network idea that the proposal put out of scope.

  \item \textbf{Val split now goes to 2020-04-30, not 2019-12-31.} \texttt{config.VAL\_END = "2020-04-30"} --- val \emph{includes} the COVID crash. The plan text and plan table both say ``2016--2019'' val. This change means the volatile-LSTM val-fallback issue (plan Phase 3b) may no longer apply, since COVID gives the val set ${\sim}2$ months of volatile-dominant windows. \textbf{Discrepancy vs plan and vs the \#2 item in the earlier discrepancy table.} The config comment on line 22 is also wrong (says test is ``2020-01-01 $\to$ present'' but test actually starts 2020-05-01).

  \item \textbf{Saved parquets are stale vs the code.} \texttt{config.LSTM\_BASELINE\_FEATURES} requires \texttt{relative\_volume\_21d}, and \texttt{features.build\_features} computes it (line 298), but the checked-in \texttt{train/val/test.parquet} files do not contain the column. Re-running notebook 01 would fix it; trying to retrain LSTMs against the committed data will \texttt{KeyError}. Also, \texttt{data/processed/test\_predictions.parquet} (which notebook 06 loads in cell 4) is not in the tree --- results only reproducible from a full local regeneration.

  \item \textbf{LSTM checkpoints are intentionally not tracked.} \texttt{.gitignore} excludes \texttt{models/*.pt}, \texttt{models/*.pth}. This is a choice (size), but combined with (4) it means a fresh clone can't reproduce the ensemble without re-running notebooks 01, 03, 04, 05, 06 in order. Worth a note in the README if we care about reproducibility for the graders.

  \item \textbf{Raw OHLCV still lives in the parquets.} Phase 6 removed them from \texttt{LSTM\_BASELINE\_FEATURES}, but the columns remain in the DataFrames (just unused). Not a bug --- feature selection happens at the config level --- but worth knowing if anyone's debugging feature plumbing.
\end{enumerate}

\hrulefill

\section*{2026-04-23 (later) --- Added Diebold--Mariano forecast significance testing}

\subsection*{What prompted this}

In discussion of the metrics table, two methodological concerns surfaced:

\begin{enumerate}[leftmargin=*]
  \item \textbf{Raw MSE is scale-misleading at this target scale.} Realized volatility sits in the 0.005--0.05 range, so MSE is necessarily tiny. Reporting ``Baseline beats Naive by 32\% MSE'' (2.2e-5 $\to$ 1.5e-5) made the gap sound dramatic, but the absolute numbers convey little about practical significance. MAPE (34.6\% $\to$ 25.3\%, a 9.3pp absolute reduction in relative error) and RMSE (0.00470 $\to$ 0.00390, ${\sim}17$\% lower) are more interpretable --- each forecast is off by roughly a quarter of the true vol, not a 30th of anything.

  \item \textbf{All rankings were point estimates.} We had not formally tested equal predictive accuracy, so we couldn't distinguish ``genuinely more accurate'' from ``slightly lower empirical loss on this sample.'' This matters most for the headline \textbf{Baseline vs Ensemble} comparison, where the RMSE gap is only ${\sim}6$\% (0.00390 vs 0.00415) --- well within the range where stochastic variation could flip the ordering.
\end{enumerate}

The second point is a legitimate weakness in the current writeup: we were making ranking claims without showing they were significant.

\subsection*{What was added}

\texttt{notebooks/06\_ensemble\_eval.ipynb} --- new \textbf{Section 6} (cells 14--17):

\begin{itemize}[leftmargin=*]
  \item \texttt{diebold\_mariano()} helper implementing the Diebold--Mariano (1995) test with the Harvey--Leybourne--Newbold (1997) small-sample correction. Newey--West HAC variance with Bartlett kernel, lag $= h - 1 = 20$ to account for autocorrelation from our overlapping 21-day targets. Supports MSE, MAE, and MAPE loss.
  \item Three comparisons, each run with both MSE and MAE loss for robustness:
  \begin{enumerate}
    \item \textbf{Naive vs Baseline LSTM} --- sanity floor.
    \item \textbf{Baseline vs Ensemble} (full test, $n=1{,}434$) --- primary research question.
    \item \textbf{Baseline vs Ensemble} on strictly-volatile days (\texttt{p\_volatile} $> 0.8$, $n=27$) --- the regime where the ensemble is supposed to help.
  \end{enumerate}
  \item A formatted results table with DM statistic, $p$-value, significance code, and $\alpha=0.05$ verdict.
  \item Interpretation notes clarifying that the strictly-volatile comparison is severely underpowered and should be read as directional evidence only.
\end{itemize}

\subsection*{Status: pending execution}

\texttt{data/processed/test\_predictions.parquet} and the \texttt{models/*.pt} LSTM checkpoints are not tracked (see ``Known reproducibility gaps''), so the DM cells cannot run on the current fresh checkout. The code is in place and will execute on the next full regeneration of notebooks 01 $\to$ 06.

\subsection*{Priors --- what we expect when it runs}

\begin{center}
\renewcommand{\arraystretch}{1.15}
\begin{tabular}{@{}p{4.8cm} p{4cm} p{5.8cm}@{}}
\toprule
Test & Prior & Reasoning \\
\midrule
Naive vs Baseline (full)            & Strongly reject ($p < 0.01$) & RMSE gap ${\sim}17$\% with $n=1{,}434$ --- ample power \\
Baseline vs Ensemble (full)         & Ambiguous                     & Gap only ${\sim}6$\% RMSE; plausibly within noise \\
Baseline vs Ensemble (volatile)     & Likely fail to reject         & $n=27$ is too small to reject despite large point-estimate gap (ensemble halves MSE) \\
\bottomrule
\end{tabular}
\end{center}

\subsection*{Why this matters for the paper}

DM results let us calibrate our claims:

\begin{itemize}[leftmargin=*]
  \item If \textbf{naive vs baseline} rejects $\to$ ``the deep model learns structure beyond persistence.'' \checkmark\ clean claim.
  \item If \textbf{baseline vs ensemble (full)} does \emph{not} reject $\to$ reframe headline from ``baseline slightly beats ensemble'' to \textbf{``baseline and ensemble are statistically tied on aggregate, and ensemble is directionally better on volatile-dominant days.''} This is both more honest and a stronger-feeling paper than claiming the ensemble lost.
  \item If \textbf{baseline vs ensemble (volatile)} does \emph{not} reject $\to$ we cannot make a significance claim on the volatile subset. Report as directional evidence plus an explicit note on low power, and suggest larger/longer test windows (extend the test period before 2020 or include additional crisis periods) as future work.
\end{itemize}

\subsection*{Weakness surfaced (for the paper's limitations section)}

Until this point we were reporting model rankings without testing them. Identified during review; now remedied for the aggregate and volatile comparisons. A remaining weakness is test-set power for regime-conditional claims: the 2020-06-30 $\to$ 2026-03-16 window contains only 27 days meeting \texttt{p\_volatile} $> 0.8$, which is insufficient to detect anything short of huge effects. This is a structural limitation of using SPY (calm-dominated) on a post-COVID test window rather than a flaw in the model.

\hrulefill

\section*{2026-04-23 (evening) --- Pipeline re-run from fresh clone; DM results executed}

\subsection*{Motivation}

Imported the repo fresh and attempted to reproduce the pipeline end-to-end. Two things we wanted: (1) confirm the reproducibility gaps we'd identified were real, (2) actually produce \texttt{test\_predictions.parquet} so the Diebold--Mariano cells added earlier today could run.

\subsection*{Reproducibility issues found (in the order we hit them)}

\begin{enumerate}[leftmargin=*]
  \item \textbf{\texttt{relative\_volume\_21d} missing from committed parquets} --- \texttt{train\_LSTM\_baseline.py} failed immediately with \texttt{KeyError: "Columns missing from dataframe: ['relative\_volume\_21d']"}. Fix: re-run nb 01. Verified this is a pre-existing gap (git status confirmed we hadn't touched \texttt{config.py}, \texttt{src/}, or \texttt{data/processed/}). Same failure any teammate would hit from a fresh clone.

  \item \textbf{\texttt{src/ensemble.py} step is undocumented} --- after nb 04 + nb 05 produced the three LSTM \texttt{.pt} files, nb 06 failed with \texttt{FileNotFoundError: data/processed/test\_predictions.parquet}. This file is produced by running \texttt{python src/ensemble.py} as an intermediate step. The README lists only notebooks; neither nb 05 nor nb 06 mentions that an external script must run between them. Fix: explicitly document the step in README, or bundle the invocation into nb 06's setup cell.

  \item \textbf{Package install was ad-hoc.} I installed \texttt{optuna}, \texttt{seaborn}, \texttt{yfinance} piecemeal as errors hit, instead of \texttt{pip install -r requirements.txt} up front. Lesson for the log: start with the README's setup block.
\end{enumerate}

\subsection*{Fresh-run metrics ($n=1{,}434$, test window 2020-06-30 $\to$ 2026-03-16)}

\begin{center}
\begin{tabular}{lrrrr}
\toprule
Model & MSE & RMSE & MAE & MAPE \\
\midrule
Naive (\texttt{rolling\_std\_21}) & 0.000022 & 0.00470 & 0.00313 & 34.6\% \\
\textbf{Baseline LSTM}            & \textbf{0.000016} & \textbf{0.00398} & \textbf{0.00257} & \textbf{25.1\%} \\
Calm LSTM                         & 0.000041 & 0.00638 & 0.00343 & 34.9\% \\
Volatile LSTM                     & 0.000085 & 0.00922 & 0.00857 & 119.4\% \\
Ensemble                          & 0.000026 & 0.00511 & 0.00310 & 31.9\% \\
\bottomrule
\end{tabular}
\end{center}

Regime-stratified ($p > 0.8$): strictly-calm $n=1{,}407$ --- Baseline MSE 0.00001, Ensemble MSE 0.00003 (baseline wins). Strictly-volatile $n=27$ --- Baseline MSE 0.00008, Ensemble MSE 0.00005 (ensemble wins directionally).

\subsection*{Diebold--Mariano results}

Sign convention: positive DM statistic means the first-listed model has higher loss (so the second-listed model is more accurate). Newey--West HAC with Bartlett kernel, lag 20; HLN small-sample correction; two-sided $t$-test with df $= n-1$.

\begin{center}
\renewcommand{\arraystretch}{1.1}
\begin{tabular}{@{}p{4cm}p{3cm}llll@{}}
\toprule
Comparison & Scope & Loss & DM & $p$ & Verdict ($\alpha=.05$) \\
\midrule
Naive vs Baseline    & full $n=1{,}434$          & MSE & $+1.36$ & 0.173 & tie \\
Naive vs Baseline    & full $n=1{,}434$          & MAE & $+2.46$ & 0.014 & \textbf{Baseline wins} \\
Baseline vs Ensemble & full $n=1{,}434$          & MSE & $-1.75$ & 0.080 & tie (marginal) \\
Baseline vs Ensemble & full $n=1{,}434$          & MAE & $-2.00$ & 0.046 & \textbf{Baseline wins} \\
Baseline vs Ensemble & volatile $n=27$           & MSE & $+0.62$ & 0.542 & tie \\
Baseline vs Ensemble & volatile $n=27$           & MAE & $+0.22$ & 0.831 & tie \\
\bottomrule
\end{tabular}
\end{center}

\subsection*{Priors vs reality}

\begin{center}
\renewcommand{\arraystretch}{1.15}
\begin{tabular}{@{}p{4cm}p{3.5cm}p{4cm}p{3.5cm}@{}}
\toprule
Test & Prior (pre-run) & Actual & Assessment \\
\midrule
Naive vs Baseline                & Strongly reject ($p < 0.01$) & MSE $p=0.17$ (tie), MAE $p=0.014$ (reject) & Weaker than predicted on MSE. The squared-error loss has higher variance, so the HAC-corrected test has less power. MAE gives the clean claim. \\
Baseline vs Ensemble (full)      & Ambiguous                     & MSE $p=0.08$, MAE $p=0.046$               & Prior confirmed. Result is borderline --- depending on loss and $\alpha$, baseline either marginally wins or ties. \\
Baseline vs Ensemble (volatile)  & Likely fail to reject despite large gap & $p=0.54$ / $0.83$                 & Prior confirmed. $n=27$ is simply too small. \\
\bottomrule
\end{tabular}
\end{center}

\subsection*{Defensible claims for the paper}

\begin{itemize}[leftmargin=*]
  \item ``The baseline LSTM is more accurate than the 21-day rolling-volatility persistence baseline (DM MAE test, $p=0.014$).'' $\leftarrow$ strongest claim; use MAE not MSE.
  \item ``The baseline LSTM matches or slightly outperforms the HMM-LSTM ensemble on aggregate (DM, $p=0.046$ MAE, $p=0.080$ MSE).'' $\leftarrow$ honest framing. Not ``the ensemble fails.''
  \item ``On volatile-dominant subperiods, the ensemble is \emph{directionally} more accurate (point-estimate MSE 0.00005 vs 0.00008), but the $n=27$ subsample is underpowered to reject equal accuracy at $\alpha=0.05$.'' $\leftarrow$ directional with explicit caveat.
\end{itemize}

\subsection*{Training stochasticity --- a new concern}

Compared to the previously-saved nb 06 output (a prior run by a teammate):

\begin{center}
\begin{tabular}{lrrl}
\toprule
Model & Prior saved MSE & This-run MSE & Change \\
\midrule
Baseline & 1.5e-5 & 1.6e-5 & +7\% \\
Calm     & 2.0e-5 & 4.1e-5 & \textbf{+105\%} \\
Volatile & 8.5e-5 & 8.5e-5 & 0\% \\
Ensemble & 1.7e-5 & 2.6e-5 & \textbf{+53\%} \\
\bottomrule
\end{tabular}
\end{center}

Same config, same seeds in \texttt{config.py} (\texttt{RANDOM\_STATE=42}, \texttt{LSTM\_RANDOM\_STATE=42}), same data. The gap comes from (a) Optuna's TPE exploration --- different trial sequences can converge to different hyperparameter neighbourhoods, (b) CPU PyTorch non-determinism without \texttt{torch.use\_deterministic\_algorithms(True)}.

\textbf{Implication:} ranking claims based on a single training run are fragile. The baseline-vs-ensemble MSE gap went from 13\% (prior) to 63\% (this run). Before locking in the narrative for the paper, we should run $k \geq 5$ seeds and report mean$\pm$std on each metric, or at minimum run 3 seeds and report the range. This is a real limitation to cite in the paper's methodology/limitations section.

\subsection*{Why this entry matters for the paper writeup}

\begin{itemize}[leftmargin=*]
  \item Shows the empirical result is genuinely nuanced, not a clean win or loss.
  \item Documents that we identified and tested methodological weaknesses (no significance testing $\to$ added DM; single-run results $\to$ flagged multi-seed need).
  \item Provides a transparent trail: priors $\to$ actual results $\to$ assessment, which is the kind of thing reviewers value.
\end{itemize}

\vspace{1em}
\hrule
\vspace{1em}

\section*{2026-04-24 --- HAR-RV benchmark design: target-formula alignment}

\subsection*{Why this entry exists}

Before coding \texttt{src/har\_rv.py} we pin down exactly what target HAR-RV will predict and exactly what predictors it will use. This is not a neutral choice --- Corsi's original HAR-RV (2002/2004) uses a different target and a different predictor convention than our project's \texttt{realized\_vol\_21d}. Getting this wrong would make the benchmark non-comparable to our LSTMs / ensembles.

User reviewed \texttt{ignore/ssrn-626064.pdf} (Corsi's original HAR-RV paper) and asked:

\begin{enumerate}[leftmargin=*]
  \item Is HAR-RV predicting backward volatility? Is that cheating?
  \item Is $\sqrt{\text{variance}} = \text{volatility}$ as assumed in our target?
  \item Document precisely where our target, Corsi's HAR-RV, and our adjusted HAR-RV agree and disagree.
\end{enumerate}

This entry answers all three.

\subsection*{Three formulas, side by side}

\paragraph{1. Our target (what every LSTM / ensemble in this project predicts).}
Defined in \texttt{src/features.py::add\_realized\_volatility()} (lines 167--216). Andersen \& Bollerslev (1998) form, adapted for daily data with demeaning:
\begin{equation}
\sigma_t \;=\; \sqrt{\,\frac{1}{m}\sum_{j=1}^{m}\bigl(r_{t+j}-\bar r\bigr)^{2}\,}
\end{equation}
with $m=21$, $\bar r$ computed over the \emph{same} forward window $[t+1,\,t+m]$.

\begin{itemize}[leftmargin=*]
  \item \textbf{Forward-looking}: aggregates returns from $t{+}1$ to $t{+}m$.
  \item \textbf{Demeaned}: AB98 subtract the window mean before squaring; appropriate for non-high-frequency data.
  \item \textbf{Normalisation $1/m$}: population-variance convention (not $1/(m{-}1)$).
  \item \textbf{Units}: volatility (sqrt of variance), same units as returns.
\end{itemize}

\paragraph{2. Corsi's original HAR-RV (ssrn-626064.pdf, equation 13).}
\begin{equation}
RV^{(d)}_{t+1d} \;=\; c \;+\; \beta^{(d)} RV^{(d)}_{t} \;+\; \beta^{(w)} RV^{(w)}_{t} \;+\; \beta^{(m)} RV^{(m)}_{t} \;+\; \omega_{t+1d}
\end{equation}
with the components defined on intraday data (his equation 3):
\begin{align}
RV^{(d)}_{t} &= \sqrt{\,\sum_{j=0}^{M-1} r^{2}_{t-j\Delta}\,}            && \text{(un-demeaned; $M$ intraday returns)} \\
RV^{(w)}_{t} &= \tfrac{1}{5}\bigl(RV^{(d)}_{t} + \ldots + RV^{(d)}_{t-4d}\bigr)  && \text{(simple 5-day average)} \\
RV^{(m)}_{t} &= \tfrac{1}{22}\bigl(RV^{(d)}_{t} + \ldots + RV^{(d)}_{t-21d}\bigr) && \text{(simple 22-day average)}
\end{align}

Structural differences from our target:

\begin{center}
\renewcommand{\arraystretch}{1.15}
\begin{tabular}{@{}p{3.2cm} p{4.8cm} p{4.8cm}@{}}
\toprule
Aspect & Our target (AB98) & Corsi HAR-RV \\
\midrule
Forecast horizon           & 21 days ahead                                    & 1 day ahead \\
Observation frequency      & daily                                            & intraday ($M$ ticks/day) \\
Demeaning                  & yes                                              & no (intraday mean $\approx 0$) \\
Normalisation              & $1/m$                                            & sum (daily) / simple avg (w,m) \\
``Monthly'' window         & 21 trading days                                  & 22 trading days \\
Units                      & volatility ($\sqrt{\cdot}$)                     & volatility ($\sqrt{\cdot}$) \\
\bottomrule
\end{tabular}
\end{center}

\textbf{Crucial agreement}: Corsi is in \emph{volatility units}, not variance. His equation (3) is $\sqrt{\cdot}$, and his regression (equation 13) is linear in those sqrt-form quantities. So $\sqrt{\text{variance}} = \text{volatility}$ is consistent between Corsi and our target. This directly answers the user's question 2.

\textbf{Not cheating}: the regression has future vol on the LHS and past-vol aggregates on the RHS. Same structure as our LSTMs --- past features $\to$ future target --- just OLS on three hand-engineered backward summaries instead of a deep model.

\paragraph{3. Our adjusted HAR-RV (what \texttt{src/har\_rv.py} will implement).}
\begin{equation}
\underbrace{\texttt{realized\_vol\_21d}(t)}_{\text{y = OUR EXACT TARGET}} \;=\; c \;+\; \beta_d \cdot RV_d(t) \;+\; \beta_w \cdot RV_w(t) \;+\; \beta_m \cdot RV_m(t) \;+\; \varepsilon
\end{equation}

with predictor components at time $t$, using past returns only:
\begin{align}
RV_d(t) &= |r_t|                                                                       && \text{(un-demeaned; mandatory, see below)} \\
RV_w(t) &= \sqrt{\,\tfrac{1}{5}\sum_{j=t-4}^{t}\bigl(r_j-\bar r_{5}\bigr)^{2}\,}       && \text{(demeaned; matches target formula)} \\
RV_m(t) &= \sqrt{\,\tfrac{1}{21}\sum_{j=t-20}^{t}\bigl(r_j-\bar r_{21}\bigr)^{2}\,}    && \text{(demeaned; matches target formula)}
\end{align}

Key invariants:

\begin{itemize}[leftmargin=*]
  \item \textbf{$y$ unchanged}: same \texttt{realized\_vol\_21d} column that every other predictor in this project (A, H, B, O, naive) is evaluated against. MSE / MAE / MAPE / Diebold--Mariano all directly comparable.
  \item \textbf{5- and 21-day components demeaned}: structurally identical to the target's formula at shorter backward windows.
  \item \textbf{1-day component un-demeaned}: forced by mathematics. Demeaning a single observation is degenerate: $\text{mean}\bigl([\,r_t\,]\bigr) = r_t$, so $(r_t - \bar r)^2 = 0$. Corsi's convention handles this the same way.
  \item \textbf{21-day monthly window (not 22)}: aligns with \texttt{config.ROLLING\_WINDOWS} and the target horizon. One-day cosmetic difference from Corsi.
\end{itemize}

\subsection*{Why the hybrid instead of pure Corsi or pure ``match-target-everywhere''}

Two pure options were considered and rejected:

\begin{itemize}[leftmargin=*]
  \item \textbf{Pure Corsi} (all predictors un-demeaned, 22-day monthly): literature-faithful but gratuitously inconsistent with our target's demeaned formula, for windows where demeaning is well-defined and numerically matters.
  \item \textbf{Pure ``match target everywhere''} (all three predictors demeaned at 1/5/21): the 1-day component mathematically collapses to 0. Benchmark reduces to a 2-predictor regression.
\end{itemize}

The hybrid preserves: $(a)$ $y$ exactly equal to our target (non-negotiable), $(b)$ predictors sharing the target's demeaned formula wherever defined (5- and 21-day), $(c)$ the 1-day component in the only sensible form available ($|r_t|$).

\subsection*{HAR-RV is a horizon-parameterised family (resolves ``doesn't HAR-RV predict tomorrow?'')}

A reasonable follow-up: Corsi's equation 13 has $RV^{(d)}_{t+1d}$ on the LHS --- i.e., \emph{tomorrow's} vol. How does that apply to our 21-day-forward problem?

``HAR-RV'' in the literature is a family parameterised by the forecast horizon $h$. The RHS (backward 1/5/22-day vol components) never changes; only the LHS target changes:
\begin{align}
h &= 1   && \to\quad RV_{t+1d}            && \text{(Corsi's original, 1-day ahead)} \\
h &= 5   && \to\quad RV_{t+1\,:\,t+5}       && \text{(weekly forecast)} \\
h &= 22  && \to\quad RV_{t+1\,:\,t+22}      && \text{(monthly forecast --- our case, $h \approx 21$)}
\end{align}

Andersen--Bollerslev--Diebold (2007), Corsi--Renò (2012), Bollerslev--Patton--Quaedvlieg (2016), and Christensen--Siggaard--Veliyev (2023) all routinely report HAR-RV at $h=1, 5, 22$ side-by-side and all call it ``HAR-RV''. The $h \approx 21$ variant is the standard monthly-horizon benchmark. Our adaptation is not exotic --- it's HAR-RV at the horizon matching our target.

\textbf{Practical expectation at $h=21$}: the monthly predictor $RV_m(t)$ will dominate the regression (large $\beta_m$), because a backward 21-day vol window is structurally the closest predictor to a forward 21-day vol target. This is itself an empirical finding worth reporting --- it's essentially why naive \texttt{rolling\_std\_21} (which is exactly the numerator of $RV_m$) is already a competitive baseline.

\subsection*{Role of the LHS: target during training vs forecast during prediction}

Second follow-up: when we write $\texttt{realized\_vol\_21d}(t) = c + \beta_d \cdot RV_d(t) + \ldots + \varepsilon$, is the LHS the \emph{actual} value or the \emph{predicted} value?

Answer: \textbf{both, depending on phase}. This is the standard regression-ML setup but worth spelling out since HAR-RV and our LSTMs share it.

\begin{center}
\renewcommand{\arraystretch}{1.25}
\begin{tabular}{@{}p{3.2cm} p{3.8cm} p{6.2cm}@{}}
\toprule
Phase & LHS is \ldots & Computed how \\
\midrule
Training (fit OLS)        & observed ground truth $y_\tau = \texttt{realized\_vol\_21d}(\tau)$ for historical $\tau \le$ train\_end & from returns $r_{\tau+1}, \ldots, r_{\tau+21}$ that have already been observed (everything in training is past, so both $X$ and $y$ are knowable) \\
Prediction (at test date $t$) & forecast $\hat y_t$                 & plug $X_t = [RV_d(t), RV_w(t), RV_m(t)]$ into the fitted equation; the true $y_t$ isn't observable until 21 trading days have elapsed \\
Evaluation (on test rows) & both available                      & compare $\hat y_t$ (forecast) to $y_t$ (realized, now computable in hindsight); MSE / MAE / DM tests run on the resulting $(\hat y_t, y_t)$ pairs \\
\bottomrule
\end{tabular}
\end{center}

\textbf{Why this is not circular}: temporal separation. $X_t$ uses only returns at times $\le t$ (past); $y_t = \texttt{realized\_vol\_21d}(t)$ uses returns at times $t{+}1, \ldots, t{+}21$ (future). At time $t$, $X_t$ is available but $y_t$ is not. The model's job is to predict $y_t$ from $X_t$. The historical training pairs $(X_\tau, y_\tau)$ exist only because training data is all in the past --- both sides are knowable \emph{to the modeller}, even though $y_\tau$ was not knowable \emph{at time $\tau$}.

\textbf{Same setup as every LSTM in this project.} The LSTMs take $X_t$ (past features) and output $\hat y_t$ (predicted forward vol). HAR-RV is $\hat y_t = \mathrm{OLS}(X_t)$ instead of $\hat y_t = \mathrm{LSTM}(X_t)$ --- same input-output contract, different function class. That's why HAR-RV's metrics are directly comparable to the LSTMs', and why the same Diebold--Mariano machinery applies to pairwise comparisons with A / H / B / O.

\subsection*{Scope notes}

\begin{itemize}[leftmargin=*]
  \item No hyperparameter tuning. HAR-RV is OLS on three features --- fit in milliseconds on $\sim$3{,}000 train rows.
  \item No log transform of the target. Our LSTMs train in log-vol space for gradient reasons; OLS has no such need and \texttt{realized\_vol\_21d} is already non-negative and well-scaled.
  \item No GARCH(1,1) companion. Scoped to HAR-RV only per user decision.
  \item The \texttt{ignore/GARCH:Review.html} file supplied for the check was a paywalled ScienceDirect abstract only --- no formulas or methodology --- so it cannot verify target alignment independently.
\end{itemize}

\subsection*{Defensible paper claim}

``We additionally compare against the HAR-RV benchmark of Corsi (2002/2004), adapted to our 21-day forward realized-volatility target: OLS regression of \texttt{realized\_vol\_21d}$(t)$ on three backward-looking volatility components at 1-, 5-, and 21-day horizons. The adaptation preserves Corsi's model structure while keeping the prediction target identical to that of our LSTM baselines, so HAR-RV's MSE / MAE are directly comparable. The 1-day component uses $|r_t|$ (un-demeaned, as in Corsi); the 5- and 21-day components use the demeaned Andersen--Bollerslev form, matching the target's own formula.''

\subsection*{Implementation pointer}

\texttt{src/har\_rv.py} --- now implemented. CLI mirrors \texttt{src/ensemble.py}: writes \texttt{data/processed/\allowbreak test\_predictions\_\allowbreak harrv.parquet} with columns \texttt{[har\_rv, target, rv\_d, rv\_w, rv\_m, naive]} indexed by date.

\subsection*{Results --- live local run (2026-04-24)}

HAR-RV fit in $\sim$100~ms (OLS on 3 features $\times$ 3{,}690 train rows). No Colab training needed for this benchmark --- HAR-RV is an econometric baseline, not a deep model.

\textbf{Learned coefficients} (match the prior: $\beta_m$ dominates at $h = 21$):

\begin{center}
\renewcommand{\arraystretch}{1.15}
\begin{tabular}{@{}l r p{7cm}@{}}
\toprule
Term & Value & Reading \\
\midrule
intercept                            & $+2.72 \times 10^{-3}$ & \\
$\beta_d$ ($|r_t|$, daily)           & $+0.046$               & tiny --- daily shock carries little 21-day signal \\
$\beta_w$ (5-day demeaned vol)       & $+0.234$               & moderate \\
$\beta_m$ (21-day demeaned vol)      & $\mathbf{+0.494}$      & \textbf{dominant} --- best-matched predictor for forward 21-day target \\
\bottomrule
\end{tabular}
\end{center}

\textbf{Test metrics} ($n = 1{,}461$; HAR-RV's backward 21-day window drops fewer leading test rows than the LSTMs' longer seq\_len):

\begin{center}
\renewcommand{\arraystretch}{1.15}
\begin{tabular}{@{}lrrrr@{}}
\toprule
Predictor & MSE & RMSE & MAE & MAPE \\
\midrule
naive (\texttt{rolling\_std\_21})        & $2.23 \times 10^{-5}$ & $0.00472$ & $0.00318$ & $34.8\%$ \\
\textbf{HAR-RV}                          & $\mathbf{1.64 \times 10^{-5}}$ & $\mathbf{0.00405}$ & $\mathbf{0.00282}$ & $\mathbf{31.5\%}$ \\
\bottomrule
\end{tabular}
\end{center}

\textbf{Diebold--Mariano vs naive} ($h = 21$, Bartlett HAC, HLN-corrected):

\begin{center}
\renewcommand{\arraystretch}{1.15}
\begin{tabular}{@{}lrrl@{}}
\toprule
Loss & DM & $p$ & Verdict ($\alpha = 0.05$) \\
\midrule
MSE          & $+1.76$ & $0.079$ & tie (marginal at $\alpha = 0.10$) \\
\textbf{MAE} & $\mathbf{+2.69}$ & $\mathbf{0.007}$ & \textbf{HAR-RV wins significantly} \\
\bottomrule
\end{tabular}
\end{center}

\subsection*{HAR-RV compared to our learned variants (informal)}

Post-sentiment-merge nb 08 numbers used the LSTM common test index ($n = 1{,}434$); HAR-RV used $n = 1{,}461$. Numbers below are \textbf{not strictly row-for-row comparable} --- proper alignment + DM pairwise HAR-RV vs each variant happens in task 6.

\begin{center}
\renewcommand{\arraystretch}{1.15}
\begin{tabular}{@{}lrrrr@{}}
\toprule
Predictor & MSE & MAE & MAPE & $n$ \\
\midrule
naive (A/H/B window)               & $2.22 \times 10^{-5}$ & $0.00313$ & $34.6\%$ & $1{,}434$ \\
naive (HAR-RV window)              & $2.23 \times 10^{-5}$ & $0.00318$ & $34.8\%$ & $1{,}461$ \\
\textbf{HAR-RV}                    & $\mathbf{1.64 \times 10^{-5}}$ & $\mathbf{0.00282}$ & $\mathbf{31.5\%}$ & $\mathbf{1{,}461}$ \\
Variant H baseline (\texttt{p\_volatile}) & $1.58 \times 10^{-5}$ & $0.00252$ & $24.9\%$ & $1{,}434$ \\
Variant A ensemble                 & $1.42 \times 10^{-5}$ & $0.00252$ & $24.6\%$ & $1{,}434$ \\
Variant B ensemble (VIX)           & $1.40 \times 10^{-5}$ & $0.00251$ & $26.9\%$ & $1{,}434$ \\
\bottomrule
\end{tabular}
\end{center}

\textbf{Interpretation (subject to row-aligned re-check in task 6):}

\begin{itemize}[leftmargin=*]
  \item HAR-RV \textbf{closes most but not all of the gap} between naive and the deep variants. MSE: naive $2.23 \times 10^{-5}$ $\to$ HAR-RV $1.64 \times 10^{-5}$ $\to$ variant A $1.42 \times 10^{-5}$. HAR-RV captures $\sim$73\% of the naive-to-A MSE improvement using just a 3-feature OLS.
  \item HAR-RV $\approx$ Variant H on MSE ($1.64 \times 10^{-5}$ vs $1.58 \times 10^{-5}$ --- $\sim$4\% gap). Variant H is a single LSTM with \texttt{p\_volatile} as a feature; HAR-RV matches it using OLS. If the DM row-aligned test returns a tie, that weakens the ``regime signal as a feature adds unique value'' argument for variant H.
  \item HAR-RV's MAPE ($31.5\%$) is noticeably worse than A/H/B ($\sim 25\%$). MAPE penalises relative errors on small-vol calm days --- HAR-RV's $\beta_m$-dominant prediction overshoots on calm days where 21-day backward vol is high but forward vol is low. The LSTMs handle this asymmetry better.
  \item \textbf{Variant A's 13-17\% MSE advantage over HAR-RV is real} in point estimate --- the deep regime-split does add something over simple backward-vol weighting. Row-aligned DM tests will confirm whether this margin is statistically significant at $\alpha = 0.05$.
\end{itemize}

\subsection*{Defensible claims for the paper (pending row-aligned re-check)}

\begin{itemize}[leftmargin=*]
  \item \emph{``The HAR-RV (Corsi 2002/2004) $h = 21$ benchmark significantly outperforms the naive 21-day persistence baseline (DM MAE $p = 0.007$), confirming that weighted multi-horizon past vol carries real signal beyond pure persistence.''} $\leftarrow$ strongest claim, directly quotable.
  \item \emph{``The regime-split LSTM ensemble (variant A) beats HAR-RV on all three point metrics, with the largest gap on MAPE (24.6\% vs 31.5\%), suggesting the deep models handle the calm-day small-vol asymmetry better than an OLS on backward RV components.''} $\leftarrow$ directional; conditional on DM confirmation.
  \item \emph{``HAR-RV is competitive with a single-LSTM regime-feature model (variant H) at $\sim$4\% MSE gap, suggesting the regime signal carries information that multi-horizon backward vol already captures.''} $\leftarrow$ conditional; strong if DM returns a tie.
\end{itemize}

\vspace{1em}
\hrule
\vspace{1em}

\section*{2026-04-24 (afternoon) --- Variant O wiring: code complete, training pending}

\subsection*{Why variant O exists}

Variant O is the \textbf{pre-everything baseline} --- trained without any of the additions we've made since Phase 6. It answers: \emph{``did any of sentiment, VIX, or the HMM-as-feature architecture actually move the needle past where we started?''}

\begin{itemize}[leftmargin=*]
  \item \textbf{HMM input}: \texttt{HMM\_PRICE\_VOLUME\_FEATURES} (11 feats --- stationary price/volume only, no sentiment).
  \item \textbf{LSTM input}: \texttt{LSTM\_STATIONARY\_FEATURES} (5 feats --- \texttt{log\_return}, \texttt{abs\_return}, \texttt{oc\_return}, \texttt{intraday\_range}, \texttt{relative\_volume\_21d}. No sentiment. No VIX. No \texttt{p\_volatile} as a feature).
  \item \textbf{Architecture}: same 3-LSTM regime-split ensemble as variant A, but with regime labels from variant O's own (sentiment-free) HMM. O and A differ \emph{only} in feature sets, not in model structure.
\end{itemize}

\subsection*{Code changes to enable variant O}

All backward-compatible --- existing A / H / B invocations keep working unchanged. Variant O activates via new CLI args.

\begin{enumerate}[leftmargin=*]
  \item \textbf{\texttt{src/train\_hmm.py}} (new module). Programmatic counterpart to \texttt{notebooks/03\_hmm\_regime.ipynb}. Parameterised on \texttt{-{}-features-group \{HMM\_PRICE\_VOLUME\_FEATURES $|$ HMM\_FEATURES\}} and \texttt{-{}-output-suffix}. Saves \texttt{hmm\_winner\{suffix\}.joblib}, \texttt{hmm\_scaler\{suffix\}.joblib}, \texttt{hmm\_meta\{suffix\}.joblib}, and \texttt{regime\_probabilities\{suffix\}.parquet}. Delegates to functions already in \texttt{src/hmm\_model.py} --- no duplication.
  \item \textbf{\texttt{src/train\_LSTM\_regime.py}}: added \texttt{-{}-regime-probs-path} and \texttt{-{}-hmm-meta-path} CLI args. Default \texttt{None} resolves to the current hardcoded paths (unchanged for A / B). Variant O passes \texttt{-{}-regime-probs-path data/processed/\allowbreak regime\_probabilities\_O.parquet} and \texttt{-{}-hmm-meta-path models/hmm\_meta\_O.joblib}. Also removed a dead \texttt{load\_viterbi()} call whose output was immediately overwritten.
  \item \textbf{\texttt{src/ensemble.py}}: same two CLI args, same default-preservation. Variant O's invocation: \texttt{python -m src.ensemble -{}-variant-suffix \_O -{}-regime-probs-path \ldots\ -{}-hmm-meta-path \ldots}.
\end{enumerate}

\subsection*{Notebook integration}

\texttt{notebooks/07\_variant\_comparison.ipynb} now has six new cells (9--14) placed after variant-B training and before ``Produce test-set predictions''. They orchestrate the variant-O pipeline end-to-end in four stages: (1) retrain HMM, (2) train baseline LSTM, (3) train regime LSTMs with variant-O regime labels, (4) run ensemble with variant-O soft probabilities.

Exact CLI invocations (see nb 08 cells 11--14 for the shell-magic variants):
\begin{small}
\begin{verbatim}
python -m src.train_hmm \
    --features-group HMM_PRICE_VOLUME_FEATURES \
    --output-suffix _O

python -m src.train_LSTM_baseline \
    --features <LSTM_STATIONARY_FEATURES> \
    --output-prefix lstm_baseline_O

python -m src.train_LSTM_regime \
    --features <LSTM_STATIONARY_FEATURES> \
    --regime both \
    --output-suffix _O \
    --regime-probs-path data/processed/regime_probabilities_O.parquet \
    --hmm-meta-path     models/hmm_meta_O.joblib

python -m src.ensemble \
    --variant-suffix _O \
    --regime-probs-path data/processed/regime_probabilities_O.parquet \
    --hmm-meta-path     models/hmm_meta_O.joblib
\end{verbatim}
\end{small}

\subsection*{Artifacts variant O will produce on first run}

\begin{itemize}[leftmargin=*]
  \item \texttt{models/hmm\_winner\_O.joblib}, \allowbreak \texttt{models/hmm\_scaler\_O.joblib}, \allowbreak \texttt{models/hmm\_meta\_O.joblib}
  \item \texttt{data/processed/regime\_probabilities\_O.parquet}
  \item \texttt{models/lstm\_baseline\_O.pt} (+ scaler), \allowbreak \texttt{models/lstm\_calm\_O.pt} (+ scaler), \allowbreak \texttt{models/lstm\_volatile\_O.pt} (+ scaler)
  \item \texttt{data/processed/test\_predictions\_O.parquet} \\ (columns: \texttt{baseline}, \texttt{calm}, \texttt{volatile}, \texttt{target}, \texttt{p\_calm}, \texttt{p\_volatile}, \texttt{ensemble})
\end{itemize}

\subsection*{Status}

\textbf{Code complete, not yet trained.} Training will run on Colab once the remaining comparison / benchmark pieces (HAR-RV module + 5-way nb-08 table cells) are in place. No training artifacts exist locally yet --- running nb 08 cells 11--14 will produce them.

\subsection*{Expected role in the 5-way comparison table}

Variant O is the \textbf{floor} row in the final comparison table $\{\text{naive}, A, H, B, O, \text{HAR-RV}\}$. Key framing questions:

\begin{itemize}[leftmargin=*]
  \item Does variant-A \textbf{ensemble with sentiment} beat variant O without sentiment? $\to$ quantifies incremental value of AAII bullish/bearish columns.
  \item Does variant H (HMM-as-feature) beat variant O? $\to$ tests whether exposing the regime signal to the LSTM via any route helps over no-regime-signal at all.
  \item Does variant B (VIX family) beat variant O? $\to$ isolates the value of forward-looking implied-vol features over price-only.
\end{itemize}

If variant O ties with A / H / B on DM tests, that's additional evidence for the structural-limits framing: \textbf{none of the post-phase-6 feature expansions mattered at this horizon}, consistent with F1 being a data-rate $\times$ horizon property rather than a feature-set problem. If variant O loses significantly, that's evidence the feature additions do carry real signal and we should reframe accordingly.

\vspace{1em}
\hrule
\vspace{1em}

\section*{2026-04-24 (evening) --- Variant-symmetric pipeline refactor}

\subsection*{Why this refactor}

Earlier pipeline had an asymmetry: nb04/nb05 trained \emph{one} ``main'' model (variant A, since \texttt{LSTM\_BASELINE\_FEATURES} defaulted to stationary + sentiment), and other variants (H, B, and the proposed O) were tacked on as ad-hoc additions in nb08. That mixed ``research question'' (\emph{does regime-splitting help?}) with ``feature-set choice'' --- confusing both for the reader and for interpretation.

The reframed research question: \textbf{``does per-regime LSTM training help over a base LSTM at predicting 21-day-forward volatility?''} --- answered \textbf{three times} on three independent pipelines (variants O / A / B), plus a separate architectural ablation (variant H).

\subsection*{Variant definitions (final, post-refactor)}

All feature lists live in \texttt{config.py}. Each variant has a matching \texttt{HMM\_VARIANT\_*\_FEATURES} (for the regime detector) and \texttt{LSTM\_VARIANT\_*\_FEATURES} (for the LSTMs).

\begin{center}
\renewcommand{\arraystretch}{1.2}
\begin{tabular}{@{}l p{3.2cm} p{3.2cm} p{4.4cm}@{}}
\toprule
Variant & HMM features & LSTM features & Research role \\
\midrule
\textbf{O} & price/vol + rolling (11) & stationary only (5) & pre-everything lower bound \\
\textbf{A} & O + sentiment (13) & O + sentiment (7) & ``main'' model (post-sentiment default) \\
\textbf{H} & \emph{(shares A's HMM)} & A's LSTM features + \texttt{p\_volatile} (8) & regime-as-feature architectural variant \\
\textbf{B} & A + VIX family (16) & A + VIX family (10) & forward-looking IV addition \\
\bottomrule
\end{tabular}
\end{center}

Backward-compat aliases preserved: \texttt{HMM\_FEATURES = HMM\_VARIANT\_A\_FEATURES}; \texttt{LSTM\_BASELINE\_FEATURES = LSTM\_VARIANT\_A\_FEATURES}; \texttt{HMM\_PRICE\_VOLUME\_FEATURES = HMM\_VARIANT\_O\_FEATURES}; \texttt{LSTM\_STATIONARY\_FEATURES = LSTM\_VARIANT\_O\_FEATURES}.

\subsection*{New notebook pipeline}

\begin{center}
\renewcommand{\arraystretch}{1.2}
\begin{tabular}{@{}l p{9.5cm}@{}}
\toprule
Notebook & Role (post-refactor) \\
\midrule
\texttt{01\_data\_collection.ipynb} & Feature engineering; new \S 7b per-variant correlation + VIF diagnostics \\
\texttt{02\_eda\_normality.ipynb}   & Normality tests extended to variant-B superset (16 features) \\
\texttt{03\_hmm\_regime.ipynb}      & Primary variant-A HMM; \S 10b trains variants O and B via \texttt{train\_hmm.py}; \S 10c injects \texttt{p\_volatile} into split parquets \\
\texttt{04\_lstm\_baseline.ipynb}   & Loops over 4 baseline LSTMs (O, A, H, B) with explicit \texttt{-{}-features} + \texttt{-{}-output-prefix} \\
\texttt{05\_lstm\_regime.ipynb}     & Loops over 3 regime pairs (O, A, B); variant-specific \texttt{-{}-regime-probs-path} + \texttt{-{}-hmm-meta-path} \\
\texttt{06\_ensemble\_eval.ipynb}   & Per-variant ensemble + \textbf{within-variant DM} (baseline vs ensemble) --- the paper's core research answer \\
\texttt{07\_variant\_comparison.ipynb} & \emph{(was \texttt{08\_seven\_model\_comparison.ipynb})} --- cross-variant 6-way comparison; 15 pairwise DM tests with Bonferroni + BH-FDR correction; F1 lag diagnostic; F2 prediction-std \\
\bottomrule
\end{tabular}
\end{center}

\subsection*{Archived notebooks}

Moved to \texttt{notebooks/archive/}:

\begin{itemize}[leftmargin=*]
  \item \texttt{07\_sentiment\_ablation.ipynb} --- superseded: the 5-feat vs 7-feat ablation is now the variant O vs A comparison in the new pipeline. Within-variant DM + cross-variant DM answer the same question more rigorously.
  \item \texttt{07\_variant\_b\_experiment.ipynb} --- already carried a deprecation notice. Pre-merge variant A was 5-feat, not the post-merge 7-feat, so its ``A vs B'' is not directly comparable to the current variant-symmetric pipeline.
\end{itemize}

\subsection*{Code changes to \texttt{src/}}

All backward-compatible --- no existing invocations break:

\begin{itemize}[leftmargin=*]
  \item \texttt{src/train\_hmm.py} (new CLI module). \texttt{-{}-features-group} accepts all three variant groups plus aliases; \texttt{-{}-output-suffix} for variant-suffixed artifacts.
  \item \texttt{src/train\_LSTM\_regime.py} --- added \texttt{-{}-regime-probs-path} and \texttt{-{}-hmm-meta-path} so each variant's regime LSTMs use its own HMM's Viterbi labels.
  \item \texttt{src/ensemble.py} --- same two CLI args for variant-specific soft-probability weighting.
  \item \texttt{config.py} --- variant-indexed feature groups; all pre-existing aliases retained.
\end{itemize}

\texttt{hmm\_model.py}, \texttt{data\_loader.py}, \texttt{lstm\_model.py}, \texttt{har\_rv.py}, \texttt{utils.py}, \texttt{train\_LSTM\_baseline.py} were not modified (already variant-agnostic by design).

\subsection*{Significance testing upgrade}

\texttt{07\_variant\_comparison.ipynb} now reports three p-value columns for each of its 30 pairwise DM tests:

\begin{itemize}[leftmargin=*]
  \item \textbf{Raw p-value} --- Diebold-Mariano (Bartlett HAC lag $h-1$, HLN correction) at $h = 21$.
  \item \textbf{Bonferroni-adjusted} --- conservative FWER at $\alpha=0.05$ requires raw $p < 0.00167$ to reject.
  \item \textbf{Benjamini-Hochberg FDR-adjusted} --- modern, less conservative standard. Controls expected proportion of false rejections at $\alpha=0.05$.
\end{itemize}

Paper claims should cite BH-FDR as the primary threshold; Bonferroni as supplementary for reviewers who expect strict FWER control.

\subsection*{Expected outcome (post-Colab)}

The research question gets \textbf{three concrete DM verdicts} in nb 06 (one per variant O / A / B), plus \textbf{15 pairwise verdicts} in nb 07 across the 6-way comparison. If the within-variant DM tests consistently return ``tie'' (baseline $\approx$ ensemble), the paper's structural-limits headline is supported three times over. If one or more reject, we quantify the effect and flag which feature-richness regime it appears in.

\emph{Note:} references to ``nb 08'' in earlier dated entries refer to what is now \texttt{07\_variant\_comparison.ipynb} after the Phase-10 rename. Historical entries are preserved as a record of the refactor trajectory.

\end{document}
